%% =============================================================================
%% section3.tex — Section 3: The Hierarchical Hurdle Beta-Binomial Model
%% =============================================================================

\section{The Hierarchical Hurdle Beta-Binomial Model}\label{sec:model}

This section develops the hierarchical hurdle beta-binomial (HBB) model
whose distributional specification was motivated by the four data features
documented in \cref{sec:features}.  The model combines three
building blocks---hurdle, beta-binomial kernel, and hierarchical
prior---into a unified framework and adds two methodological
contributions that are new to this class of models: a cross-margin
covariance whose identification operates entirely through the
hierarchical prior rather than the conditional likelihood, and a
Cholesky-based sandwich correction that extends
pseudo-posterior inference to the two-part hierarchical setting.

In plain terms, the model works as follows.  A \emph{hurdle} first
separates providers that serve infants and toddlers from those that do
not, treating participation as a logistic regression with
state-specific intercepts and slopes.  For participants, a
\emph{beta-binomial} kernel models the share of enrollment devoted to
infants and toddlers, accommodating the bounded count and within-state
overdispersion.  The two margins share a joint hierarchical prior on
their state-level deviations, so that a state's participation barrier
and its conditional enrollment intensity can be positively or negatively
correlated---the cross-margin covariance that is the paper's central
structural novelty.  Sampling weights enter through a
pseudo-likelihood, and a sandwich correction restores design-consistent
coverage for population-average parameters.

%% ---------------------------------------------------------------------------
%% 3.1 Notation and Data Structure
%% ---------------------------------------------------------------------------

\subsection{Notation and data structure}\label{sec:notation}

We observe $N = 6{,}785$ center-based childcare providers drawn from
$S = 51$ states (including the District of Columbia).  Each provider $i$
is nested in a state $s[i]$.  The outcome pair $(Y_i, n_i)$ records
infant/toddler enrollment $Y_i \in \{0, 1, \ldots, n_i\}$ out of total
enrollment $n_i$; the participation indicator is $z_i = \one(Y_i > 0)$.
\cref{tab:notation} collects the symbols used throughout the paper.
All continuous covariates are standardized prior to analysis.

\begin{table}[ht]
\centering
\caption{Notation summary. Dimensions are shown in parentheses.}\label{tab:notation}
\smallskip
\begin{tabularx}{\textwidth}{@{}lclX@{}}
\toprule
Symbol & Dimension & Domain & Description \\
\midrule
\multicolumn{4}{@{}l}{\textit{Indices and dimensions}} \\[2pt]
$i$ & --- & $\{1,\ldots,N\}$ & Provider index ($N = 6{,}785$) \\
$s$ & --- & $\{1,\ldots,S\}$ & State index ($S = 51$, incl.\ DC) \\
$s[i]$ & --- & $\{1,\ldots,S\}$ & State membership of provider $i$ \\
$N_s$ & --- & $\mathbb{Z}_{>0}$ & Number of providers in state $s$; $\sum_s N_s = N$ \\[4pt]
\multicolumn{4}{@{}l}{\textit{Outcome variables}} \\[2pt]
$n_i$ & $(1)$ & $\mathbb{Z}_{>0}$ & Total enrollment, ages 0--5 (treated as fixed) \\
$Y_i$ & $(1)$ & $\{0,1,\ldots,n_i\}$ & Infant/toddler (IT) enrollment count \\
$z_i$ & $(1)$ & $\{0,1\}$ & $z_i = \one(Y_i > 0)$: IT participation indicator \\[4pt]
\multicolumn{4}{@{}l}{\textit{Covariates}} \\[2pt]
$\bx_i$ & $(P \times 1)$ & $\R^P$ & Provider covariates, $P=5$: intercept, poverty, urban, Black, Hispanic \\
$\bx_i^{(r)}$ & $(q\times 1)$ & $\R^q$ & State-varying subvector of $\bx_i$ ($q = P$ in M2/M3) \\
$q_i$ & $(1)$ & $(0,1)$ & Participation probability (extensive margin); cf.~\eqref{eq:eta-ext} \\
$\bv_s$ & $(Q\times 1)$ & $\R^Q$ & State policy vector, $Q = 4$: intercept + $Q-1$ policy instruments \\[4pt]
\multicolumn{4}{@{}l}{\textit{Parameters}} \\[2pt]
$\balpha$ & $(P \times 1)$ & $\R^P$ & Population-average extensive-margin coefficients \\
$\bbeta$ & $(P \times 1)$ & $\R^P$ & Population-average intensive-margin coefficients \\
$\kappa$ & $(1)$ & $\R_{>0}$ & Beta-binomial concentration \\
$\bdelta_{k,s}$ & $(q\times 1)$ & $\R^q$ & State deviations, margin $k\in\{1,2\}$, state $s$ \\
$\bGamma_k$ & $(q\times Q)$ & $\R^{q \times Q}$ & Cross-level policy moderator matrix, margin $k$ \\
$\bSigma_\delta$ & $(2q\times 2q)$ & $\mathbb{S}_{++}^{2q}$ & Cross-margin state-effect covariance \\[4pt]
\multicolumn{4}{@{}l}{\textit{Survey design}} \\[2pt]
$w_i$ & $(1)$ & $\R_{>0}$ & Sampling weight \\
$\tilde{w}_i$ & $(1)$ & $\R_{>0}$ & Normalized weight: $\tilde{w}_i = w_i N / \textstyle\sum_j w_j$ \\
\bottomrule
\end{tabularx}
\end{table}

%% ---------------------------------------------------------------------------
%% 3.2 The Hurdle Beta-Binomial Likelihood
%% ---------------------------------------------------------------------------

\subsection{The hurdle beta-binomial likelihood}\label{sec:likelihood}

\paragraph{The beta-binomial distribution.}
The beta-binomial distribution in the mean--precision parameterization
\citep{Griffiths1973,Prentice1986}
$(\mu,\kappa)$ sets $a = \mu\kappa$ and $b = (1-\mu)\kappa$, where
$\mu \in (0,1)$ is the conditional mean parameter and $\kappa > 0$ is the
concentration parameter. Larger $\kappa$ implies less overdispersion, with the
$\Bin(n,\mu)$ distribution recovered as $\kappa \to \infty$. The probability
mass function (PMF) is
\begin{equation}\label{eq:bb-pmf}
  \Prob(Y = y \mid n,\mu,\kappa)
  = \binom{n}{y}
    \frac{B\bigl(y + \mu\kappa,\; n - y + (1-\mu)\kappa\bigr)}
         {B\bigl(\mu\kappa,\;(1-\mu)\kappa\bigr)},
  \quad y = 0,1,\ldots,n,
\end{equation}
where $B(\cdot,\cdot)$ is the beta function. The first two moments are
\begin{equation}\label{eq:bb-moments}
  \E[Y] = n\mu, \qquad
  \Var(Y) = n\mu(1-\mu)\,\frac{n + \kappa}{1 + \kappa}.
\end{equation}
The variance inflation factor $(n+\kappa)/(1+\kappa)$ exceeds unity for all
$n \ge 2$, capturing overdispersion relative to the binomial. In the
childcare data, with a typical enrollment of $n_i \approx 50$ and an
estimated $\kappa \approx 7$, this factor is approximately $57/8 \approx 7$,
far beyond what a binomial model can accommodate.

\paragraph{Computational remark.}
In practice, we compute $\log f_{\BetaBin}$ using the
\texttt{lbeta()} function to avoid numerical overflow when $n_i$
exceeds 50. Specifically, the log-PMF is evaluated as
$\log\binom{n}{y} + \texttt{lbeta}(y+a,\,n-y+b) - \texttt{lbeta}(a,b)$,
which remains stable for $n_i$ up to several hundred.

\paragraph{Zero probability.}
At $y = 0$ the PMF simplifies to a product form
that is both interpretable and numerically stable:
\begin{equation}\label{eq:p0}
  p_0(n,\mu,\kappa)
  = \prod_{j=0}^{n-1} \frac{(1-\mu)\kappa + j}{\kappa + j}.
\end{equation}
Each factor in the product is the conditional probability that the $j$-th
``trial'' yields a non-IT outcome, given all preceding trials did so. The
product form shows that $p_0$ is strictly decreasing in both $\mu$ and $\kappa$
(for $n \ge 2$): higher mean intensity and lower overdispersion each reduce the
probability of observing zero IT enrollment.

\paragraph{Hurdle construction.}
As documented in \cref{sec:features}, the zeros in the childcare data are
structural: they reflect a deliberate decision not to serve infants rather
than a random draw from a count distribution.  This motivates a hurdle
specification \citep{Mullahy1986,Cragg1971} that separates the participation
decision from the enrollment intensity.

\begin{definition}[Hurdle beta-binomial model]\label{def:hbb}
  For provider $i$ with enrollment capacity $n_i$, the HBB model specifies
  \begin{equation}\label{eq:hbb}
    f(y_i \mid q_i, \mu_i, \kappa, n_i) =
    \begin{cases}
      1 - q_i, & \text{if } y_i = 0, \\[4pt]
      q_i \cdot f_{\ZTBB}(y_i \mid n_i, \mu_i, \kappa),
        & \text{if } y_i \in \{1,\ldots,n_i\},
    \end{cases}
  \end{equation}
  where $q_i \in (0,1)$ is the participation probability (extensive margin),
  $\mu_i \in (0,1)$ is the conditional mean IT share (intensive margin), and
  the zero-truncated beta-binomial PMF is
  \begin{equation}\label{eq:ztbb}
    f_{\ZTBB}(y \mid n,\mu,\kappa)
    = \frac{f_{\BetaBin}(y \mid n,\mu,\kappa)}{1 - p_0(n,\mu,\kappa)},
    \quad y = 1,\ldots,n.
  \end{equation}
\end{definition}

The hurdle separates two distinct decisions that childcare centers face. The
\emph{extensive margin} captures whether a center serves any infants at
all---the access decision---governed by $q_i$. The \emph{intensive margin}
captures how many infant slots a center offers relative to its total
capacity---the intensity decision---governed by $\mu_i$. This separation
reflects a genuine two-stage process: a center must first decide to accept the
costs and regulatory burden of infant care (clearing the hurdle), and only then
does it determine how many infant slots to offer.

The log-likelihood under conditional independence decomposes additively into
extensive and intensive components,
\begin{multline}\label{eq:loglik-decomp}
  \ell(\btheta)
  = \underbrace{\sum_{i=1}^{N}\bigl[z_i\log q_i + (1-z_i)\log(1-q_i)\bigr]}
    _{\ell_{\mathrm{ext}}(\balpha)} \\
  + \underbrace{\sum_{i:\,z_i=1}\bigl[\log f_{\BetaBin}(y_i \mid n_i,\mu_i,\kappa)
    - \log\bigl(1 - p_{0,i}\bigr)\bigr]}
    _{\ell_{\mathrm{int}}(\bbeta,\kappa)},
\end{multline}
where $p_{0,i} = p_0(n_i,\mu_i,\kappa)$. This separation is a structural
property of hurdle models: the binary indicator $z_i$ is a sufficient statistic
partition that renders the extensive parameters $\balpha$ and the intensive
parameters $(\bbeta,\kappa)$ variation-independent in the likelihood. In
contrast, zero-inflated models introduce latent at-risk indicators that couple
the two components, complicating both estimation and interpretation.

\paragraph{Unconditional expectation and the intensity function.}
The unconditional mean enrollment count is
\begin{equation}\label{eq:uncond-mean}
  \E[Y_i] = q_i \cdot n_i \cdot h(\mu_i, n_i, \kappa),
  \qquad
  h(\mu, n, \kappa) \equiv \frac{\mu}{1 - p_0(n,\mu,\kappa)},
\end{equation}
where $h$ is the \emph{intensity function} mapping the latent mean proportion
$\mu$ to the expected proportion among participants. Since $h(\mu) \ge \mu$
(truncation shifts the conditional mean upward), participating centers are
predicted to have a higher IT share than the latent parameter $\mu$ alone
would suggest. The function $h$ is central to the reversal
interpretation: the key question is whether the opposing signs of the extensive
and intensive coefficients can be ``absorbed'' by a nonlinear response of $h$,
or whether the reversal is genuinely a coefficient-level phenomenon. The
following theorem resolves this.

\begin{theorem}[Monotonicity of $h$]\label{thm:monotonicity}
  For all $n \ge 2$, $\kappa > 0$, and $\mu \in (0,1)$,
  \begin{equation}\label{eq:dh-dmu}
    \frac{\partial h}{\partial \mu} > 0.
  \end{equation}
  For $n = 1$, $h(\mu) \equiv 1$ for all $\mu \in (0,1)$.
\end{theorem}

\begin{proof}[Proof sketch]
  Define $\Phi(\mu) = (1 - p_0) - \mu\, p_0\, \Lambda$, where
  $\Lambda = \kappa\sum_{j=0}^{n-1}[(1-\mu)\kappa + j]^{-1} > 0$.
  Since $\partial h / \partial\mu = \Phi(\mu)/(1-p_0)^2$, the sign of the
  derivative equals the sign of $\Phi$. One verifies that $\Phi(0) = 0$ and
  $\Phi'(\mu) = \mu\, p_0(\Lambda^2 - \Lambda_2)$, where
  $\Lambda_2 = \kappa^2\sum_{j=0}^{n-1}[(1-\mu)\kappa + j]^{-2}$.
  For $n \ge 2$, the Cauchy--Schwarz inequality gives
  $\Lambda^2 > \Lambda_2$, so $\Phi'(\mu) > 0$ on $(0,1)$.
  Since $\Phi$ starts at zero and is strictly increasing, $\Phi(\mu) > 0$ for
  all $\mu \in (0,1)$. The full proof appears in~\cref{app:math-properties}.
\end{proof}

\paragraph{Interpretation for the poverty reversal.}
\cref{thm:monotonicity} establishes that the intensity function $h$ is strictly
increasing in $\mu$: a center with a higher underlying IT intensity $\mu$ will
always have a higher expected IT share among participants, regardless of how much
overdispersion is present. Consequently, the poverty reversal---whereby higher poverty reduces participation ($\alpha_{\mathrm{pov}} < 0$) yet
increases conditional intensity ($\beta_{\mathrm{pov}} > 0$)---is purely a
\emph{coefficient-level} phenomenon, not an artifact of the nonlinear response
function. Researchers can examine the posterior distributions of
$\alpha_{\mathrm{pov}}$ and $\beta_{\mathrm{pov}}$ directly to assess the
reversal, without needing to account for nonlinear link function effects.

\begin{proposition}[Elasticity of $h$]\label{prop:elasticity}
  Define the elasticity $\varepsilon_h = (\partial h / \partial\mu)(\mu / h)$.
  Then
  \begin{equation}\label{eq:elasticity}
    \varepsilon_h(\mu) = 1 - \omega(\mu),
    \qquad
    \omega(\mu) = \frac{\mu\, p_0\, \Lambda}{1 - p_0} \in (0,1)
    \quad \text{for } n \ge 2.
  \end{equation}
  Hence $\varepsilon_h \in (0,1)$: the intensity function responds
  inelastically to changes in $\mu$.
\end{proposition}

The inelastic response arises because increasing $\mu$ simultaneously increases
the probability that a center would have had a positive draw even without the
hurdle (raising $1 - p_0$, the denominator of $h$), partially offsetting the
numerator effect. In practical terms, a one-percent increase in a center's
underlying IT propensity $\mu$ translates into \emph{less} than a one-percent
increase in its expected IT share among participants. At our estimated parameters
($\mu \approx 0.30$, $\kappa \approx 7$, $n \approx 50$), $\omega \approx 0.07$,
so the elasticity is approximately 0.93---relatively modest damping in this
parameter region.

%% ---------------------------------------------------------------------------
%% 3.3 State-Varying Coefficients with Cross-Margin Covariance
%% ---------------------------------------------------------------------------

\subsection{State-varying coefficients (SVC) with cross-margin covariance}\label{sec:svc}

The reversal is not a single national pattern. Preliminary estimates
show that fewer than half of states exhibit the classic reversal---a negative
extensive-margin poverty coefficient paired with a positive intensive-margin
one (\cref{sec:prelim}). This geographic variation raises a fundamental
question: \emph{Does the poverty--enrollment relationship vary systematically
across states, and if so, are the extensive and intensive margins linked?}

\paragraph{Linear predictors.}
Following the spatially varying coefficient framework of~\citet{GelfandEtAl2003},
the extensive and intensive margins are linked to provider characteristics
through logit linear predictors:
\begin{align}
  \logit(q_i)
    &= \bx_i\t\balpha + \bx_i^{(r)\t}\bdelta_{1,s[i]},
    \label{eq:eta-ext} \\[4pt]
  \logit(\mu_i)
    &= \bx_i\t\bbeta + \bx_i^{(r)\t}\bdelta_{2,s[i]},
    \label{eq:eta-int}
\end{align}
where $\balpha, \bbeta \in \R^P$ are population-average fixed effects---the
``national story''---and $\bdelta_{k,s} \in \R^q$
($k = 1$ for extensive, $k = 2$ for intensive) are zero-mean state-specific
deviations that capture how each state's poverty--enrollment relationship
departs from the national average. The subvector $\bx_i^{(r)}$ consists of
the first $q$ components of $\bx_i$ whose coefficients are permitted to vary
across states.
The \emph{total state-specific coefficient} for covariate $j$ on margin $k$ is
\begin{equation}\label{eq:total-coef}
  \tilde{\alpha}_{j,s} = \alpha_j + \delta_{1,s,j},
  \qquad
  \tilde{\beta}_{j,s} = \beta_j + \delta_{2,s,j}.
\end{equation}
In the childcare application, a state with $\tilde{\alpha}_{\mathrm{pov},s} < 0$
and $\tilde{\beta}_{\mathrm{pov},s} > 0$ exhibits the reversal: higher
community poverty is associated with lower participation but higher conditional
intensity.

\paragraph{Cross-margin covariance: the key novelty.}
Standard hurdle models estimate the two margins independently. We depart from
this convention by stacking the state deviations into a single $2q$-dimensional
vector:
\begin{equation}\label{eq:delta-stack}
  \bdelta_s
  = \begin{pmatrix} \bdelta_{1,s} \\ \bdelta_{2,s} \end{pmatrix}
  \sim \Norm_{2q}\!\parens{\mathbf{0},\;\bSigma_\delta},
  \qquad s = 1,\ldots,S,
\end{equation}
with the $2q \times 2q$ positive-definite covariance
\begin{equation}\label{eq:sigma-delta}
  \bSigma_\delta
  = \begin{pmatrix}
      \bSigma_{11} & \bSigma_{12} \\
      \bSigma_{21} & \bSigma_{22}
    \end{pmatrix},
\end{equation}
where $\bSigma_{11}$ governs between-state variation in extensive-margin
coefficients, $\bSigma_{22}$ governs the intensive margin, and the
off-diagonal block $\bSigma_{12} = \bSigma_{21}\t$ captures
\emph{cross-margin dependence} (in the policy-moderated model M3b, $\bSigma_{12}$ refers to the \emph{residual} cross-margin covariance after removing policy effects; see~\cref{smb:rem-sigma12-residual}). This cross-margin covariance is the central
structural novelty of the model: it allows states with stronger participation
barriers (larger $|\delta_{1,\mathrm{pov},s}|$) to simultaneously exhibit
stronger or weaker intensity effects ($\delta_{2,\mathrm{pov},s}$), with the
direction and magnitude estimated from data. The cross-margin correlation for
covariate $j$ is
\begin{equation}\label{eq:cross-corr}
  \varrho_j^{\mathrm{cross}} = \frac{[\bSigma_{12}]_{jj}}
  {\sqrt{[\bSigma_{11}]_{jj}\,[\bSigma_{22}]_{jj}}}.
\end{equation}

\begin{proposition}[Non-identification of $\bSigma_{12}$ from the
  conditional likelihood]
  \label{prop:sigma12-nonid}
  Conditional on the state-level deviations $\bdelta_s$, the Fisher
  information matrix for the HBB model is block-diagonal between
  the extensive-margin parameters $(\balpha, \bdelta_{1,1},\ldots,\bdelta_{1,S})$
  and the intensive-margin parameters
  $(\bbeta, \bdelta_{2,1},\ldots,\bdelta_{2,S}, \kappa)$. Consequently,
  $\bSigma_{12}$ receives zero information from the conditional
  likelihood and is identified through the hierarchical prior on
  $\bdelta_s$ acting on the $S = 51$ estimated state-level deviations.
\end{proposition}

\begin{proof}
  The log-likelihood decomposes as $\ell = \ell_{\mathrm{ext}} + \ell_{\mathrm{int}}$
  (\cref{eq:loglik-decomp}). Since $\ell_{\mathrm{ext}}$ depends only on
  $(\balpha, \bdelta_{1,\cdot})$ and $\ell_{\mathrm{int}}$ depends only on
  $(\bbeta, \bdelta_{2,\cdot}, \kappa)$, the cross-derivatives
  $\partial^2 \ell / \partial\phi_1 \partial\phi_2\t$ vanish for any
  extensive-margin parameter $\phi_1$ and intensive-margin parameter $\phi_2$.
  Hence the observed and expected information matrices are block-diagonal
  across margins, and the off-diagonal block $\bSigma_{12}$ appears nowhere
  in the likelihood.
\end{proof}

\paragraph{$\bSigma_{12}$ as a hierarchical estimand.}
\cref{prop:sigma12-nonid} clarifies the information structure for
$\bSigma_{12}$: its posterior is shaped entirely by the $S = 51$ pairs
$(\bdelta_{1,s}, \bdelta_{2,s})$ drawn from the common
$\Norm_{2q}(\mathbf{0}, \bSigma_\delta)$ distribution, giving it an
effective sample size of $S$ rather than $N$.  While a frequentist could
in principle estimate $\bSigma_{12}$ via REML on the marginal
likelihood, the high-dimensional covariance structure ($2q = 10$
free parameters per margin) relative to the number of groups ($S = 51$)
makes unrestricted frequentist estimation fragile.  The $\LKJ(\eta)$
prior provides finite-sample regularization, concentrating the posterior
at rate $O(S^{-1/2})$. This contrasts with the within-margin blocks
$\bSigma_{11}$ and $\bSigma_{22}$, which also receive direct likelihood
information through the state-specific coefficients and concentrate at
rate $O(N^{-1/2})$.

\begin{remark}[Conditional versus marginal likelihood]\label{rem:cond-vs-marg}
  The block-diagonality in \cref{prop:sigma12-nonid} is a statement
  about the \emph{conditional} likelihood $p(\mathbf{y} \mid \bdelta,
  \balpha, \bbeta, \kappa)$.  The marginal likelihood obtained by
  integrating over the random effects,
  $p(\mathbf{y} \mid \bSigma_\delta, \balpha, \bbeta, \kappa) = \int
  p(\mathbf{y} \mid \bdelta) \, p(\bdelta \mid \bSigma_\delta) \,
  \mathrm{d}\bdelta$, does \emph{not} factorize when $\bSigma_{12}
  \neq \mathbf{0}$, so that $\bSigma_{12}$ is in principle identified
  through the marginal likelihood.  However, the conditional
  formulation is natural for MCMC computation, where $\bdelta_s$ is
  sampled as a latent variable and $\bSigma_\delta$ is updated
  conditional on the realized deviations.  In a frequentist framework,
  REML estimation of $\bSigma_{12}$ via the marginal likelihood is
  possible in principle but numerically fragile when $S = 51$ groups
  contribute to a $2q = 10$-dimensional covariance matrix: boundary
  solutions and singular Hessians are common in this regime.  The
  Bayesian approach with $\LKJ$ regularization provides stable
  estimation at the cost of the slower $O(S^{-1/2})$ concentration
  rate documented in \cref{prop:sigma12-id}.
\end{remark}

\begin{proposition}[Identifiability of $\bSigma_{12}$ via the prior]
  \label{prop:sigma12-id}
  Suppose (i) $S > 2q$, (ii) each state $s$ has at least $q$ providers with
  linearly independent covariates, and (iii) the prior on $\bSigma_\delta$ has
  full support on $\mathbb{S}_{++}^{2q}$. Then the posterior of $\bSigma_\delta$
  (including $\bSigma_{12}$) is well-concentrated, with the posterior standard
  deviation of each element of $\bSigma_{12}$ decreasing at rate $O(S^{-1/2})$.
  For our data, $S = 51 \gg 2q = 10$.
\end{proposition}

\begin{remark}[Boundary condition]\label{rem:boundary}
  When $S < 2q$ the cross-margin covariance is poorly identified even with a
  proper prior, because $S$ draws from a $2q$-dimensional distribution yield a
  singular sample covariance matrix. In this regime, the scalar cross-margin
  parameter $a$ of~\citet{GhosalEtAl2020}, which nests as
  $\bSigma_{12} = a\,\bSigma_{22}$, provides a parsimonious but still useful
  summary of cross-margin dependence.
\end{remark}

\paragraph{Computational remark.}
In Stan~\citep{CarpenterEtAl2017}, the state deviations are sampled via the non-centered
parameterization (NCP): draw $\bz_s \iid \Norm_{2q}(\mathbf{0},\mathbf{I})$
and set $\bdelta_s = \diag(\boldsymbol{\tau})\,\mathbf{L}_\varepsilon\,\bz_s$,
where $\mathbf{L}_\varepsilon$ is the Cholesky factor of the
correlation matrix $\mathbf{R}_\varepsilon$. This eliminates the funnel
geometry~\citep{BetancourtGirolami2015} that arises under centered parameterization when small $\tau_j$ forces
$\delta_{s,j}$ near zero, causing divergent transitions in Hamiltonian Monte
Carlo. With $S = 51$ states and $N_s$ ranging from 17 to 568, NCP is the
appropriate default.

%% ---------------------------------------------------------------------------
%% 3.4 Policy Moderators and Marginal Effect Decomposition
%% ---------------------------------------------------------------------------

\subsection{Cross-level policy moderators and marginal effect decomposition}
\label{sec:policy}

Why does the reversal appear in some states but not others?  State-level
subsidy policies provide a natural explanation: they directly affect the cost
structure that centers face when deciding whether and how much to serve infants.
We formalize this hypothesis through cross-level policy moderators.

\paragraph{The moderation framework.}
\begin{definition}[Policy moderator specification]\label{def:gamma}
  For each margin $k \in \{1,2\}$, the state-specific deviation decomposes as
  \begin{equation}\label{eq:gamma-decomp}
    \bdelta_{k,s} = \bGamma_k\,\bv_s + \bepsilon_{k,s},
  \end{equation}
  where $\bGamma_k \in \R^{q \times Q}$ is the cross-level interaction matrix
  and $\bepsilon_{k,s}$ is a residual state effect.
  Stacking both margins, the residuals follow a joint distribution that
  preserves cross-margin dependence:
  \begin{equation}\label{eq:eps-stacked}
    \bepsilon_s
    \equiv \begin{pmatrix}\bepsilon_{1,s}\\\bepsilon_{2,s}\end{pmatrix}
    \sim \Norm_{2q}\!\left(\mathbf{0},\;
      \bSigma_\varepsilon
      = \begin{pmatrix}
        \bSigma_{\varepsilon,1} & \bSigma_{12}\\
        \bSigma_{21} & \bSigma_{\varepsilon,2}
      \end{pmatrix}\right),
  \end{equation}
  so that the full model estimates the residual cross-margin covariance
  $\bSigma_{12} = \Cov(\bepsilon_{1,s},\, \bepsilon_{2,s})$
  after removing the policy-explained component (see
  \cref{prop:sigma12-nonid} and Supplementary Material~B).
\end{definition}

The state-level policy vector is
$\bv_s \in \R^Q$ with $Q = 4$, consisting of an intercept and $Q - 1 = 3$
Child Care and Development Fund (CCDF) policy instruments that directly affect
the cost structure centers face for infant care.  The specific policy variables
and their substantive roles are detailed in \cref{sec:application}.

\paragraph{Nested model hierarchy.}
The HBB framework nests a sequence of models with increasing complexity,
summarized in~\cref{tab:model-hierarchy}.

\begin{table}[ht]
\centering
\caption{Nested model hierarchy. Each model adds structure relative to its
  predecessor.}\label{tab:model-hierarchy}
\smallskip
\small
\begin{tabular}{@{}llp{0.38\textwidth}c@{}}
\toprule
Model & Structure & Question addressed & Approx.\ params \\
\midrule
M0 & Fixed effects only & National-level poverty reversal? & $11$ \\
M1 & + Random intercepts$^{\dagger}$ & Does enrollment vary by state? & $116$ \\
M2 & + State-varying coefs.\ & Does the reversal vary by state? & $\sim 576$ \\
M3a & + Cross-margin cov.\ & Are the two margins linked? & $\sim 576$ \\
M3b & + Policy moderators & Can policies explain state variation? & $\sim 626$ \\
\bottomrule
\end{tabular}
\par\vspace{2pt}
{\footnotesize $^{\dagger}$M1 is obtained as a special case of M2 by
  restricting the SVC to intercepts only ($q = 1$).}
\end{table}

Each transition addresses a question raised by its predecessor: M0 to M1 asks
whether the baseline enrollment rate varies by state; M1 to M2, whether the
reversal pattern varies by state; M2 to M3a, whether cross-margin
covariance improves fit; and M3a to M3b, whether observable policies explain
the state heterogeneity.

\paragraph{Poverty coefficient expansion.}
\begin{proposition}[Poverty coefficient expansion]\label{prop:poverty-expansion}
  Under M3b, the total state-specific poverty coefficient on the extensive
  margin is
  \begin{equation}\label{eq:poverty-expansion}
    \tilde{\alpha}_{\mathrm{pov},s}
    = \underbrace{\alpha_{\mathrm{pov}}}_{\substack{\text{national} \\
    \text{average}}}
    + \underbrace{\bgamma_{1,\mathrm{pov}}\t\bv_s}_{\substack{\text{policy-explained}
    \\ \text{deviation}}}
    + \underbrace{\varepsilon_{1,\mathrm{pov},s}}_{\substack{\text{residual}
    \\ \text{state effect}}},
    %% Note: \varepsilon_{k,\cdot,s} denotes residual state deviations
    %% (distinct from the elasticity \varepsilon_h in Proposition~\ref{prop:elasticity}).
  \end{equation}
  where $\bgamma_{1,\mathrm{pov}} = (\gamma_{1,2,1},\ldots,\gamma_{1,2,Q})\t$
  is the second row of $\bGamma_1$. Analogously, the intensive-margin
  counterpart is
  $\tilde{\beta}_{\mathrm{pov},s} = \beta_{\mathrm{pov}} +
  \bgamma_{2,\mathrm{pov}}\t\bv_s + \varepsilon_{2,\mathrm{pov},s}$.
\end{proposition}

\cref{prop:poverty-expansion} decomposes the state-specific poverty effect into
three interpretable components: a national average ($\alpha_{\mathrm{pov}}$), a
policy-explained deviation ($\bgamma_{1,\mathrm{pov}}\t\bv_s$), and a residual
idiosyncrasy ($\varepsilon_{1,\mathrm{pov},s}$).  The variance of the
state-specific poverty coefficient decomposes correspondingly as
$\Var(\tilde{\alpha}_{\mathrm{pov},s}) =
\bgamma_{1,\mathrm{pov}}\t\,\Var(\bv_s)\,\bgamma_{1,\mathrm{pov}}
+ (\bSigma_{\varepsilon,1})_{\mathrm{pov,pov}}$,
partitioning between-state variation into policy-explained and residual
components.

\paragraph{Marginal effect decomposition.}
To translate model parameters into policy-relevant quantities, we decompose the
total effect of any covariate $x_k$ on expected enrollment into access and
intensity channels.

\begin{proposition}[Log-access and log-intensity decomposition]
  \label{prop:lae-lie}
  For provider $i$ in state $s$, the marginal effect of covariate $x_k$ on
  log expected enrollment decomposes as
  \begin{equation}\label{eq:lae-lie}
    \frac{\partial \log \E[Y_i]}{\partial x_{i,k}}
    = \underbrace{(1 - q_i)\,\tilde{\alpha}_{k,s}}_{\LAE_k}
    \;+\;
    \underbrace{(1 - \mu_i)\,\varepsilon_{h,i}\,\tilde{\beta}_{k,s}}_{\LIE_k},
  \end{equation}
  where $\LAE_k$ is the log access effect and $\LIE_k$ is the log
  intensity effect, with $\varepsilon_{h,i} = 1 - \omega_i$ from
  \cref{prop:elasticity}.
\end{proposition}

\begin{proof}
  Since $\log\E[Y_i] = \log q_i + \log n_i + \log h_i$ and $n_i$ is fixed,
  the chain rule gives two terms. For the access term, the logit link yields
  $\partial(\log q_i)/\partial x_{i,k} = (1 - q_i)\,\tilde{\alpha}_{k,s}$.
  For the intensity term,
  $\partial(\log h_i)/\partial x_{i,k}
  = \varepsilon_{h,i}\,(1-\mu_i)\,\tilde{\beta}_{k,s}$,
  using the logistic derivative $\partial\mu/\partial\eta = \mu(1-\mu)$
  and the elasticity definition
  $\varepsilon_{h} = \mu\,\partial(\log h)/\partial\mu$.
\end{proof}

\paragraph{Concrete interpretation.}
At the sample average ($q \approx 0.64$, $\mu \approx 0.30$,
$\varepsilon_h \approx 0.93$) with population-average coefficients
$\alpha_{\mathrm{pov}} = -0.324$ and $\beta_{\mathrm{pov}} = +0.090$:
$\LAE_{\mathrm{pov}} = (1 - 0.64)(-0.324) = -0.117$ and
$\LIE_{\mathrm{pov}} = (1 - 0.30)(0.93)(0.090) = +0.058$.
The net effect is $-5.9\%$: the access barrier dominates, but the intensity
response offsets nearly half.  Three features of the decomposition merit
emphasis: it is \emph{exactly additive} on the log scale, both multipliers
$(1-q_i)$ and $(1-\mu_i)\varepsilon_{h,i}$ are positive and bounded so that
the sign of each component is determined solely by the coefficient sign, and
the reversal pattern is confirmed as a coefficient-level phenomenon consistent
with~\cref{thm:monotonicity}.

%% ---------------------------------------------------------------------------
%% 3.5 Survey-Weighted Pseudo-Posterior and Sandwich Correction
%% ---------------------------------------------------------------------------

\subsection{Survey-weighted pseudo-posterior and sandwich correction}
\label{sec:survey}

The NSECE 2019 employs a stratified multistage cluster design with 30 strata,
415 primary sampling units (PSUs), and sampling weights ranging from 1 to 462
(Kish $\ESS = 1{,}803$; $\DEFF = 3.76$). Under informative sampling designs,
ignoring these weights can bias point estimates and produce
miscalibrated credible intervals
\citep{PfeffermannEtAl1998}. Following~\citet{SavitskyToth2016} and
\citet{WilliamsSavitsky2021}, we adopt a pseudo-posterior approach with
post-hoc sandwich correction
\citep[see also][for a general treatment of weighting in multilevel
models]{RabeHeskethSkrondal2006}.

\paragraph{Pseudo-log-likelihood.}
Define the weighted log-likelihood
\begin{equation}\label{eq:pseudo-loglik}
  \ell^{(w)}(\btheta)
  = \sum_{i=1}^{N} \tilde{w}_i \log f_{\HBB}(y_i \mid \btheta),
\end{equation}
where $\tilde{w}_i = w_i N / \sum_j w_j$ are the normalized survey weights.

\begin{theorem}[Pseudo-posterior propriety]\label{thm:propriety}
  Under a proper prior $\pi(\btheta)$, the pseudo-posterior
  \begin{equation}\label{eq:pseudo-posterior}
    \pi^{(w)}(\btheta \mid \by)
    \propto \exp\!\bigl(\ell^{(w)}(\btheta)\bigr)\,\pi(\btheta)
  \end{equation}
  is proper for any set of positive, finite weights $\{w_i\}_{i=1}^N$.
\end{theorem}

\begin{proof}[Proof sketch]
  Since $f_{\HBB}(y_i \mid \btheta) \in (0,1]$ and $\tilde{w}_i > 0$,
  each weighted factor satisfies $f_i^{\tilde{w}_i} \le 1$. Hence
  $\int \exp(\ell^{(w)})\,\pi(\btheta)\,d\btheta
  \le \int \pi(\btheta)\,d\btheta = 1 < \infty$, using
  the properness of $\pi(\btheta)$. The full proof, which verifies that the
  marginal likelihood is strictly positive, appears in Supplementary
  Material~C.
\end{proof}

\begin{proposition}[Weighted separability]\label{prop:weighted-sep}
  The weighted log-likelihood retains the two-part decomposition:
  \begin{equation}\label{eq:weighted-sep}
    \ell^{(w)}(\btheta)
    = \ell_{\mathrm{ext}}^{(w)}(\balpha)
    + \ell_{\mathrm{int}}^{(w)}(\bbeta, \kappa),
  \end{equation}
  where
  $\ell_{\mathrm{ext}}^{(w)} = \sum_i \tilde{w}_i[z_i\log q_i + (1-z_i)\log(1-q_i)]$
  and
  $\ell_{\mathrm{int}}^{(w)} = \sum_{i:\,z_i=1}
  \tilde{w}_i[\log f_{\BetaBin}(y_i \mid n_i,\mu_i,\kappa) - \log(1-p_{0,i})]$.
\end{proposition}

\begin{proof}
  Weighting by $\tilde{w}_i > 0$ multiplies each observation's log-density
  by a positive scalar, which cannot introduce functional dependence between
  parameters that were already in separate additive components.
\end{proof}

\paragraph{Bernstein--von Mises (BvM) condition.}
The weight ratio $w_{\max}/w_{\min} = 462$ exceeds $N^{1/2} \approx 82.5$.
The finite-sample diagnostic
$\delta \approx \log(462)/\log(6{,}785) \approx 0.695 > 0.5$ suggests that
the pseudo-posterior variance may understate the true sampling variability,
motivating the sandwich correction below.

\paragraph{Sandwich variance.}
Define the observed Hessian and cluster-robust outer product:
\begin{align}
  \mathbf{H}_{\mathrm{obs}}
    &= -\sum_{i=1}^{N}\tilde{w}_i\,
       \frac{\partial^2\log f_{\HBB}(y_i \mid \btheta)}
            {\partial\btheta\,\partial\btheta\t}
       \bigg|_{\btheta = \hat{\btheta}},
    \label{eq:H-obs} \\[6pt]
  \mathbf{J}_{\mathrm{cluster}}
    &= \sum_{c=1}^{C}\mathbf{s}_c\,\mathbf{s}_c\t,
    \qquad
    \mathbf{s}_c
    = \sum_{i \in \text{PSU}_c}\tilde{w}_i\,
      \frac{\partial\log f_{\HBB}(y_i \mid \btheta)}{\partial\btheta}
      \bigg|_{\btheta = \hat{\btheta}},
    \label{eq:J-cluster}
\end{align}
where $C = 415$ is the number of primary sampling units and $\hat{\btheta}$ is
the posterior mean. The hurdle separability implies that
$\mathbf{H}_{\mathrm{obs}}$ is block-diagonal between the extensive and
intensive margin parameters; however, $\mathbf{J}_{\mathrm{cluster}}$ is
\emph{not} block-diagonal, because a single PSU contributes scores from both
margins simultaneously.

\begin{remark}[Simplified exposition]\label{rem:J-simplified}
  \Cref{eq:J-cluster} uses the unstratified form for expositional
  clarity.  The implemented estimator is the stratified cluster-robust
  form in~\cref{smc:prop-cluster-robust}, which centers PSU scores
  within the $H = 30$ NSECE strata and applies finite-population
  correction factors $C_h/(C_h - 1)$.
\end{remark}

The design-consistent sandwich variance is
\begin{equation}\label{eq:V-sand}
  \mathbf{V}_{\mathrm{sand}}
  = \mathbf{H}_{\mathrm{obs}}^{-1}\,
    \mathbf{J}_{\mathrm{cluster}}\,
    \mathbf{H}_{\mathrm{obs}}^{-1}.
\end{equation}

\begin{theorem}[Cholesky affine transformation]\label{thm:cholesky}
  Let $\hat{\btheta}$ be the posterior mean, $\mathbf{L}_{\mathrm{sand}}$ the
  lower Cholesky factor of $\mathbf{V}_{\mathrm{sand}}$, and
  $\mathbf{L}_{\mathrm{MCMC}}$ the lower Cholesky factor of the MCMC posterior
  covariance $\bSigma_{\mathrm{MCMC}}$. The transformed draws
  \begin{equation}\label{eq:cholesky-transform}
    \btheta^{*(m)}
    = \hat{\btheta}
    + \mathbf{L}_{\mathrm{sand}}\,\mathbf{L}_{\mathrm{MCMC}}^{-1}
      \bigl(\btheta^{(m)} - \hat{\btheta}\bigr),
    \qquad m = 1,\ldots,M,
  \end{equation}
  satisfy $\E[\btheta^*] = \hat{\btheta}$ (mean-preserving) and
  $\Cov(\btheta^*) = \mathbf{V}_{\mathrm{sand}}$
  (design-consistent covariance).  The transformation matches
  the target covariance structure; Wald confidence intervals
  based on the marginal variances are invariant to the choice
  of transformation matrix.
\end{theorem}

\paragraph{Computational remark.}
The observed Hessian $\mathbf{H}_{\mathrm{obs}}$ is computed as the negative
second derivative of the weighted pseudo-log-likelihood~\eqref{eq:H-obs},
excluding prior curvature, evaluated at the posterior mean,
\emph{not} from the MCMC sample covariance. This distinction is critical in
hierarchical models: the marginal posterior variance
$\bSigma_{\mathrm{MCMC}}$ of the fixed effects integrates over the
random-effect distribution, producing intervals much wider than the
conditional precision $\mathbf{H}_{\mathrm{obs}}^{-1}$. In our
application, the hierarchical variance ratio
$(\bSigma_{\mathrm{MCMC}})_{pp}/(\mathbf{H}_{\mathrm{obs}}^{-1})_{pp}$
exceeds 600 for most fixed effects---a consequence of the hierarchical
structure, not prior domination---confirming that the standard substitution
$\bSigma_{\mathrm{MCMC}} \approx \mathbf{H}^{-1}$ fails in hierarchical
models.

\paragraph{Design effect ratio and block-wise correction strategy.}
As a parameter-specific diagnostic extending the classical design
effect~\citep{RaoScott1981, PfeffermannEtAl1998} to individual model
parameters, define the design effect ratio for parameter $p$ as
\begin{equation}\label{eq:der}
  \DER_{p}
  = \frac{[\mathbf{V}_{\mathrm{sand}}]_{pp}}{[\mathbf{H}_{\mathrm{obs}}^{-1}]_{pp}}.
\end{equation}
The $\DER_p$ quantifies the factor by which the survey design inflates
the variance of the $p$-th parameter beyond the data-only (SRS) variance
and motivates the following three-tier correction strategy.
The Cholesky transformation~\eqref{eq:cholesky-transform} is applied
in blocks:
\begin{itemize}[nosep]
  \item \emph{Fixed effects} $(\balpha, \bbeta, \log\kappa)$: full sandwich
    correction. The DER ranges from 1.14 to 4.18 in our data
    (mean 2.11, consistent with the Kish $\DEFF$ of 3.76). The corresponding
    confidence interval inflation factors are
    $\sqrt{\DER} \in [1.07, 2.04]$.
  \item \emph{Random effects} $(\bdelta_{k,s})$: partial correction guided by
    the DER. States with low effective sample size show DER near one (the prior
    dominates), while large states show DER up to four.
  \item \emph{Hyperparameters} $(\bSigma_\delta, \bGamma_k)$: no correction
    applied. The sandwich estimator is not applicable to variance components
    estimated from between-state variation; the DER concept does not extend to
    these parameters. Consequently, posterior intervals for hyperparameters are
    model-based summaries whose validity relies on correct specification of the
    hierarchical distribution, not on design consistency. When a hyperparameter
    is substantively central---as $\bSigma_{12}$ is in this
    application (\cref{prop:sigma12-nonid,prop:sigma12-id})---prior
    sensitivity analysis is recommended.
\end{itemize}

\begin{remark}[Block-wise caveat]\label{rem:blockwise}
  The block-wise application of~\cref{thm:cholesky} ignores cross-block
  covariance between fixed and random effects, which can be non-negligible
  for the three to five largest states where within-state ESS approaches $N_s$.
  A sensitivity analysis comparing block-wise and joint correction is reported
  in Supplementary Material~C.  The finite-sample coverage of sandwich-corrected
  intervals is evaluated in the simulation study (\cref{sec:simulation}).
\end{remark}

%% ---------------------------------------------------------------------------
%% 3.6 Prior Specification and Computation
%% ---------------------------------------------------------------------------

\subsection{Prior specification and computation}\label{sec:priors}

\paragraph{Prior hierarchy.}
Each prior is chosen to be weakly informative---providing enough
regularization to ensure computational stability and proper posteriors, while
remaining vague enough that the data dominate in well-identified regions.
\cref{tab:priors} summarizes the complete specification, following the
recommendations of~\citet{GelmanEtAl2008} for regularized regression and
\citet{Gelman2006} for variance components.

\begin{table}[ht]
\centering
\caption{Prior specification. NCP = non-centered parameterization.}\label{tab:priors}
\smallskip
\footnotesize
\begin{adjustbox}{max width=\textwidth}
\begin{tabular}{@{}llp{0.50\textwidth}@{}}
\toprule
Parameter & Prior & Justification \\
\midrule
$\balpha,\;\bbeta$ & $\Norm(\mathbf{0},\;2^2\,\mathbf{I}_P)$ &
  Weakly informative on logit scale; \\
  & & $\pm 4$ logit units covers most of the probability range \\[3pt]
$\operatorname{vec}(\bGamma_k)$ & $\Norm(0,\;1^2)$ &
  Diffuse; allows moderate cross-level effects \\[3pt]
$\log\kappa$ & $\Norm(2,\;1.5^2)$ &
  Centered at moderate overdispersion ($\kappa \approx 7.4$); \\
  & & 95\% prior interval $\kappa \in [0.4, 140]$ \\[3pt]
$\boldsymbol{\tau}$ (SDs of $\bSigma_\delta$) &
  $\HalfNormal(0,1)$ &
  Regularizing; shrinkage-friendly \\[3pt]
$\mathbf{R}_\varepsilon$ (correlations) & $\LKJ(2)$ \citep{LewandowskiEtAl2009} &
  Mild shrinkage toward identity \\[3pt]
$\bz_s$ (NCP auxiliaries) & $\Norm(\mathbf{0},\;\mathbf{I}_{2q})$ &
  Standard normal (non-centered parameterization) \\
\bottomrule
\end{tabular}
\end{adjustbox}
\end{table}

\paragraph{Prior for $\log\kappa$.}
We center the log-concentration prior at $2$ (corresponding to $\kappa \approx 7.4$)
because preliminary analysis of IT enrollment counts suggests moderate
overdispersion: the variance-to-mean ratio across participating centers is
approximately 7--15 times the binomial baseline, consistent with $\kappa$ in the
range 5--20.
The scale of 1.5 is generous, accommodating $\kappa$ values from less than~1
(near-maximum overdispersion) to over~140 (approaching the binomial).

\paragraph{Identification.}
\begin{theorem}[Identification]\label{thm:identification}
  Suppose the following conditions hold:
  \begin{enumerate}[label=\textup{(C\arabic*)},nosep]
    \item The design matrices have full rank:
      $\mathrm{rank}(\mathbf{X}) = P$ and $\mathrm{rank}(\mathbf{V}) = Q$.
    \item Each state $s$ has at least $q$ providers with linearly independent
      covariate vectors $\bx_i^{(r)}$, with at least one zero and at least
      $q$ positive responders.
    \item $n_i \ge 3$ for all $i$ with $z_i = 1$, and there exist providers
      $i, j$ with $s[i] = s[j]$ and $n_i \ne n_j$ in each state.
    \item The prior $\pi(\btheta)$ has full support on the parameter space.
  \end{enumerate}
  Then the full parameter vector $\btheta = (\balpha, \bbeta, \kappa,
  \{\bdelta_s\}, \bSigma_\delta, \{\bGamma_k\})$ is identified under the
  posterior.
\end{theorem}

\noindent\emph{Interpretation.}
The hurdle likelihood identifies the state-specific \emph{total}
coefficients $\tilde{\alpha}_{j,s}$ and $\tilde{\beta}_{j,s}$
(\cref{eq:total-coef}) under standard rank and replication
conditions; the hierarchical layer then pins down their
decomposition into population-average effects~$\balpha^{(r)}$,
$\bbeta^{(r)}$, policy moderation~$\bGamma_k$, and residual state
deviations~$\bepsilon_{k,s}$.
See Supplementary Material~B for the complete proof.

Condition (C3) is critical for the intensive margin: with $n_i = 1$, the
intensity function $h \equiv 1$ and the parameters $\mu_i$ and $\kappa$
enter the likelihood only through $1 - p_0 = \mu_i$, so $\kappa$ is not
identified. With $n_i = 2$, $p_0$ depends on both $\mu$ and $\kappa$, but the
single constraint $p_0(2,\mu,\kappa) = (1-\mu)[(1-\mu)\kappa+1]/(\kappa+1)$ determines
only one degree of freedom. For $n_i \ge 3$, the product formula
\eqref{eq:p0} provides enough structure to separate $\mu$ from $\kappa$. The
full proof appears in Supplementary Material~B.

\paragraph{Estimation pipeline.}
\cref{alg:estimation} outlines the complete estimation and inference procedure.

\begin{algorithm}[ht]
\caption{Estimation pipeline for the HBB model}\label{alg:estimation}
\begin{algorithmic}[1]
  \State \textbf{Data preparation:} Normalize weights
    $\tilde{w}_i = w_i N / \textstyle\sum_j w_j$; standardize covariates.
  \State \textbf{Compile:} Stan model with NCP and \texttt{lbeta()} for numerical
    stability.
  \State \textbf{MCMC sampling:} Run 4 chains~\citep[NUTS;][]{Hoffman2014}
    with \texttt{adapt\_delta} $= 0.95$; see Supplementary Material~E for
    iteration counts and tree-depth settings.
  \State \textbf{Convergence diagnostics:} Verify $\hat{R} < 1.01$,
    $\ESS_{\mathrm{bulk}} > 400$~\citep{Vehtari2021}, and zero divergent transitions for all
    parameters.
  \State \textbf{Score extraction:} Extract observation-level scores
    $\partial\log f_{\HBB} / \partial\btheta$ computed in Stan's
    \texttt{generated quantities} block, evaluated at each MCMC draw.
  \State \textbf{Sandwich computation:} Form $\mathbf{H}_{\mathrm{obs}}$
    and $\mathbf{J}_{\mathrm{cluster}}$ via~\cref{eq:H-obs,eq:J-cluster};
    compute $\mathbf{V}_{\mathrm{sand}}$ via~\cref{eq:V-sand}.
  \State \textbf{Cholesky transformation:} Apply
    \cref{eq:cholesky-transform} block-wise to obtain design-corrected
    posterior draws $\{\btheta^{*(m)}\}$.
  \State \textbf{Model comparison:} Compute LOO-CV
    \citep{VehtariEtAl2017} across M0--M3b; report
    $\widehat{\elpd}_{\mathrm{loo}}$ and Pareto-$k$ diagnostics.
\end{algorithmic}
\end{algorithm}

The models are fitted sequentially from M0 through M3b, using the posterior
mean of each simpler model as the starting point for its successor.
This warm-start strategy reduces total computation time by approximately
40\% compared with cold starts.  Full MCMC settings and convergence diagnostics
are reported in Supplementary Material~E; score computation details and
closed-form expressions for the sandwich ingredients are documented in the
\textsf{hurdlebb} package vignette~\citep{hurdlebb2026}.

The finite-sample properties of this estimation pipeline---particularly
the sandwich correction's behavior under realistic survey designs---are
evaluated in the simulation study that follows.
